\documentclass[11pt,a4paper]{article}
\usepackage[margin=2.2cm]{geometry}
\usepackage{xcolor}
\usepackage{titlesec}
\usepackage{enumitem}
\usepackage{mdframed}
\usepackage{fancyhdr}
\usepackage{hyperref}
\usepackage{cmap}
\usepackage[T1]{fontenc}
\usepackage{parskip}
\usepackage{booktabs}
\usepackage{tabularx}
\usepackage{amssymb}
\usepackage{fancyvrb}
\usepackage{fvextra}

% ── Colours ──────────────────────────────────────────────────────────────────
\definecolor{primary}{RGB}{25,70,150}
\definecolor{accent}{RGB}{0,150,120}
\definecolor{warn}{RGB}{190,70,20}
\definecolor{purple}{RGB}{110,50,160}
\definecolor{codebg}{RGB}{244,246,250}
\definecolor{boxbg}{RGB}{232,242,255}
\definecolor{boxborder}{RGB}{25,70,150}
\definecolor{dangercol}{RGB}{255,238,235}
\definecolor{dangerborder}{RGB}{190,50,20}
\definecolor{taskcol}{RGB}{245,240,255}
\definecolor{taskborder}{RGB}{110,50,160}
\definecolor{notecol}{RGB}{255,250,225}
\definecolor{noteborder}{RGB}{200,160,0}
\definecolor{treecol}{RGB}{245,245,255}
\definecolor{treeborder}{RGB}{110,50,160}
\definecolor{greencol}{RGB}{235,252,240}
\definecolor{greenborder}{RGB}{0,140,60}

% ── Code block: plain verbatim, fully copyable ───────────────────────────────
% codeblock environment: framed box with light background, monospace font
% Uses fancyvrb's BVerbatim so it can sit inside mdframed
\newmdenv[
  backgroundcolor=codebg,
  linecolor=gray!40,
  linewidth=0.6pt,
  leftline=true,rightline=true,topline=true,bottomline=true,
  innerleftmargin=10pt,innerrightmargin=10pt,
  innertopmargin=7pt,innerbottommargin=7pt,
  splittopskip=0pt
]{codeframe}

% Convenience command: \begin{codeblock} ... \end{codeblock}
% wraps fvextra Verbatim inside a codeframe box
\DefineVerbatimEnvironment{codeblock}{Verbatim}{
  fontsize=\footnotesize,
  fontfamily=tt,
  breaklines=true,
  breakanywhere=false,
  breaksymbolleft={},
  tabsize=4
}

% ── Sections ─────────────────────────────────────────────────────────────────
\titleformat{\section}{\large\bfseries\color{primary}}{}{0em}{}[\color{primary}\titlerule]
\titleformat{\subsection}{\normalsize\bfseries\color{accent}}{}{0em}{\textbullet\enspace}
\titleformat{\subsubsection}{\normalsize\bfseries\color{purple}}{}{0em}{\small$\triangleright$\enspace}

% ── Info boxes ───────────────────────────────────────────────────────────────
\newmdenv[backgroundcolor=boxbg,linecolor=boxborder,linewidth=2.5pt,
  leftline=true,rightline=false,topline=false,bottomline=false,
  innerleftmargin=10pt,innerrightmargin=10pt,
  innertopmargin=7pt,innerbottommargin=7pt]{infobox}

\newmdenv[backgroundcolor=dangercol,linecolor=dangerborder,linewidth=2.5pt,
  leftline=true,rightline=false,topline=false,bottomline=false,
  innerleftmargin=10pt,innerrightmargin=10pt,
  innertopmargin=7pt,innerbottommargin=7pt]{dangerbox}

\newmdenv[backgroundcolor=taskcol,linecolor=taskborder,linewidth=2.5pt,
  leftline=true,rightline=false,topline=false,bottomline=false,
  innerleftmargin=10pt,innerrightmargin=10pt,
  innertopmargin=7pt,innerbottommargin=7pt]{taskbox}

\newmdenv[backgroundcolor=notecol,linecolor=noteborder,linewidth=2.5pt,
  leftline=true,rightline=false,topline=false,bottomline=false,
  innerleftmargin=10pt,innerrightmargin=10pt,
  innertopmargin=7pt,innerbottommargin=7pt]{notebox}

\newmdenv[backgroundcolor=treecol,linecolor=treeborder,linewidth=2pt,
  leftline=true,rightline=false,topline=false,bottomline=false,
  innerleftmargin=10pt,innerrightmargin=10pt,
  innertopmargin=7pt,innerbottommargin=7pt]{treebox}

\newmdenv[backgroundcolor=greencol,linecolor=greenborder,linewidth=2.5pt,
  leftline=true,rightline=false,topline=false,bottomline=false,
  innerleftmargin=10pt,innerrightmargin=10pt,
  innertopmargin=7pt,innerbottommargin=7pt]{greenbox}

% ── Caption label for code blocks ────────────────────────────────────────────
\newcommand{\codelabel}[1]{%
  \vspace{2pt}%
  {\footnotesize\color{gray}\textit{#1}}%
  \vspace{2pt}%
}

% ── Header / Footer ──────────────────────────────────────────────────────────
\pagestyle{fancy}
\fancyhf{}
\fancyhead[L]{\small\color{primary}\textbf{C\# Interview Training --- ShopFlow}}
\fancyhead[R]{\small\color{gray}Stage 0 --- Project Setup \& Broken Scaffold}
\fancyfoot[C]{\small\color{gray}\thepage}
\renewcommand{\headrulewidth}{0.4pt}
\hypersetup{colorlinks=true,linkcolor=primary,urlcolor=accent}

% =============================================================================
\begin{document}

% ── Title page ────────────────────────────────────────────────────────────────
\begin{titlepage}
\centering
\vspace*{2.5cm}
{\Huge\bfseries\color{primary}ShopFlow\par}
\vspace{0.25cm}
{\large\color{gray}Order Processing \& Inventory Management API\par}
\vspace{0.9cm}
\rule{0.6\textwidth}{1.5pt}
\vspace{0.9cm}
{\LARGE\bfseries\color{warn}Stage 0\par}
\vspace{0.2cm}
{\large Project Setup, Architecture \& Broken Scaffold\par}
\vspace{1.8cm}

\begin{infobox}
\textbf{What Stage 0 is}\\[6pt]
You will scaffold the full solution from scratch using the terminal,
create every project and wire up references, then be handed a large chunk
of broken production code spread across multiple files.\\[6pt]
Your job is to \textbf{identify every problem, fix it, and write tests}.
Paste your solution back in the chat.\\[6pt]
Your coach will then critique it, show the ideal version, and set Stage 1.
\end{infobox}

\vspace{1.2cm}
\begin{center}
\begin{tabular}{ll}
\textbf{Domain:}      & E-commerce order processing \& inventory \\
\textbf{Stack:}       & ASP.NET Core 8, EF Core, xUnit, Moq \\
\textbf{Deliverable:} & Fixed code pasted in chat \\
\textbf{Tests:}       & Minimum 14 unit tests \\
\end{tabular}
\end{center}

\vfill
{\small\color{gray}March 2026 --- C\# Senior Developer Interview Training}
\end{titlepage}

\tableofcontents
\newpage

% =============================================================================
\section{Project Overview --- ShopFlow}

ShopFlow is a backend API for a mid-size e-commerce company.
It manages the full order lifecycle: customers browse products,
place orders, make payments, and receive notifications.
Warehouse staff manage inventory and fulfil shipments.

By the end of all stages the system will have:

\begin{itemize}[leftmargin=1.5em,itemsep=3pt]
  \item A fully layered Clean Architecture solution (Domain, Application,
        Infrastructure, API)
  \item A rich domain model: Orders, Products, Inventory, Customers, Payments
  \item Repository pattern and Unit of Work backed by EF Core
  \item REST API with proper HTTP semantics and global error handling
  \item Unit tests and integration tests with full isolation
  \item Background jobs for stock reservation expiry
  \item Docker Compose for local development
\end{itemize}

\begin{notebox}
\textbf{Stage 0 scope:} Scaffold the solution and fix the broken service layer.
No EF Core, no HTTP endpoints yet --- those come in later stages.
\end{notebox}

% =============================================================================
\newpage
\section{Solution Setup --- Run These Commands}

Run every block in order. Each block is one commit.

% ─────────────────────────────────────────────────────────────────────────────
\subsection{Step 1 --- Create the Solution Root}

\begin{codeframe}
\begin{codeblock}
mkdir ShopFlow && cd ShopFlow
dotnet new sln --name ShopFlow
git init
git commit --allow-empty -m "chore: init repo"
\end{codeblock}
\end{codeframe}

% ─────────────────────────────────────────────────────────────────────────────
\subsection{Step 2 --- Create .gitignore}

Create a file named \texttt{.gitignore} in the root:

\begin{codeframe}
\begin{codeblock}
bin/
obj/
.vs/
*.user
.env
*.db
.idea/
\end{codeblock}
\end{codeframe}

\begin{codeframe}
\begin{codeblock}
git add .gitignore
git commit -m "chore: add gitignore"
\end{codeblock}
\end{codeframe}

% ─────────────────────────────────────────────────────────────────────────────
\subsection{Step 3 --- Scaffold All Projects}

\begin{codeframe}
\begin{codeblock}
# Core layers
dotnet new classlib -n ShopFlow.Domain         -o src/ShopFlow.Domain
dotnet new classlib -n ShopFlow.Application    -o src/ShopFlow.Application
dotnet new classlib -n ShopFlow.Infrastructure -o src/ShopFlow.Infrastructure
dotnet new webapi   -n ShopFlow.Api            -o src/ShopFlow.Api

# Test projects
dotnet new xunit -n ShopFlow.Tests.Unit        -o tests/ShopFlow.Tests.Unit
dotnet new xunit -n ShopFlow.Tests.Integration -o tests/ShopFlow.Tests.Integration
\end{codeblock}
\end{codeframe}

% ─────────────────────────────────────────────────────────────────────────────
\subsection{Step 4 --- Register Projects in the Solution}

\begin{codeframe}
\begin{codeblock}
dotnet sln add src/ShopFlow.Domain/ShopFlow.Domain.csproj
dotnet sln add src/ShopFlow.Application/ShopFlow.Application.csproj
dotnet sln add src/ShopFlow.Infrastructure/ShopFlow.Infrastructure.csproj
dotnet sln add src/ShopFlow.Api/ShopFlow.Api.csproj
dotnet sln add tests/ShopFlow.Tests.Unit/ShopFlow.Tests.Unit.csproj
dotnet sln add tests/ShopFlow.Tests.Integration/ShopFlow.Tests.Integration.csproj
\end{codeblock}
\end{codeframe}

% ─────────────────────────────────────────────────────────────────────────────
\subsection{Step 5 --- Wire Up Project References}

\begin{codeframe}
\begin{codeblock}
# Application depends on Domain
dotnet add src/ShopFlow.Application    reference src/ShopFlow.Domain

# Infrastructure implements Application interfaces
dotnet add src/ShopFlow.Infrastructure reference src/ShopFlow.Application

# Api is the composition root
dotnet add src/ShopFlow.Api            reference src/ShopFlow.Application
dotnet add src/ShopFlow.Api            reference src/ShopFlow.Infrastructure

# Unit tests cover Domain + Application
dotnet add tests/ShopFlow.Tests.Unit   reference src/ShopFlow.Domain
dotnet add tests/ShopFlow.Tests.Unit   reference src/ShopFlow.Application

# Integration tests spin up the Api host
dotnet add tests/ShopFlow.Tests.Integration reference src/ShopFlow.Api
\end{codeblock}
\end{codeframe}

% ─────────────────────────────────────────────────────────────────────────────
\subsection{Step 6 --- Add NuGet Packages}

\begin{codeframe}
\begin{codeblock}
# Unit test packages
dotnet add tests/ShopFlow.Tests.Unit package xunit
dotnet add tests/ShopFlow.Tests.Unit package xunit.runner.visualstudio
dotnet add tests/ShopFlow.Tests.Unit package Moq
dotnet add tests/ShopFlow.Tests.Unit package FluentAssertions
dotnet add tests/ShopFlow.Tests.Unit package coverlet.collector

# Integration test packages (used from Stage 3 onwards)
dotnet add tests/ShopFlow.Tests.Integration package Microsoft.AspNetCore.Mvc.Testing
dotnet add tests/ShopFlow.Tests.Integration package FluentAssertions

git add .
git commit -m "chore: scaffold solution, projects, references and packages"
\end{codeblock}
\end{codeframe}

% ─────────────────────────────────────────────────────────────────────────────
\subsection{Step 7 --- Create the Internal Folder Structure}

\begin{codeframe}
\begin{codeblock}
# Domain
mkdir -p src/ShopFlow.Domain/Entities
mkdir -p src/ShopFlow.Domain/ValueObjects
mkdir -p src/ShopFlow.Domain/Enums
mkdir -p src/ShopFlow.Domain/Exceptions

# Application
mkdir -p src/ShopFlow.Application/Interfaces
mkdir -p src/ShopFlow.Application/Services
mkdir -p src/ShopFlow.Application/DTOs
mkdir -p src/ShopFlow.Application/Common

# Infrastructure (stubbed -- implemented in later stages)
mkdir -p src/ShopFlow.Infrastructure/Persistence
mkdir -p src/ShopFlow.Infrastructure/Repositories
mkdir -p src/ShopFlow.Infrastructure/Services
mkdir -p src/ShopFlow.Infrastructure/Extensions

# Api (stubbed -- implemented in later stages)
mkdir -p src/ShopFlow.Api/Controllers
mkdir -p src/ShopFlow.Api/Middleware
mkdir -p src/ShopFlow.Api/Extensions

# Test folders
mkdir -p tests/ShopFlow.Tests.Unit/Domain
mkdir -p tests/ShopFlow.Tests.Unit/Application
mkdir -p tests/ShopFlow.Tests.Unit/Helpers
mkdir -p tests/ShopFlow.Tests.Integration/Endpoints

git add .
git commit -m "chore: create internal folder structure"
\end{codeblock}
\end{codeframe}

% =============================================================================
\newpage
\section{Full Directory Tree}

After all steps above your solution must look exactly like this:

\begin{treebox}
\begin{codeblock}
ShopFlow/
|-- ShopFlow.sln
|-- .gitignore
|
|-- src/
|   |
|   |-- ShopFlow.Domain/
|   |   |-- ShopFlow.Domain.csproj        <- zero NuGet refs allowed
|   |   |-- Entities/
|   |   |   |-- Order.cs
|   |   |   |-- OrderItem.cs
|   |   |   |-- Product.cs
|   |   |   |-- Customer.cs
|   |   |   `-- InventoryEntry.cs
|   |   |-- ValueObjects/
|   |   |   |-- Money.cs
|   |   |   `-- Address.cs
|   |   |-- Enums/
|   |   |   |-- OrderStatus.cs
|   |   |   `-- PaymentStatus.cs
|   |   `-- Exceptions/
|   |       |-- DomainException.cs
|   |       `-- NotFoundException.cs
|   |
|   |-- ShopFlow.Application/
|   |   |-- ShopFlow.Application.csproj
|   |   |-- Interfaces/
|   |   |   |-- IOrderRepository.cs
|   |   |   |-- IProductRepository.cs
|   |   |   |-- IInventoryRepository.cs
|   |   |   |-- IEmailSender.cs
|   |   |   `-- IUnitOfWork.cs
|   |   |-- Services/
|   |   |   |-- OrderService.cs           <- fix this
|   |   |   |-- InventoryService.cs       <- fix this
|   |   |   `-- PromoEngine.cs            <- create this
|   |   |-- DTOs/
|   |   |   |-- PlaceOrderRequest.cs
|   |   |   |-- OrderDto.cs
|   |   |   `-- OrderItemDto.cs
|   |   `-- Common/
|   |       `-- Result.cs
|   |
|   |-- ShopFlow.Infrastructure/
|   |   |-- ShopFlow.Infrastructure.csproj
|   |   |-- Repositories/
|   |   |   |-- InMemoryOrderRepository.cs   <- implement this
|   |   |   `-- InMemoryProductRepository.cs <- implement this
|   |   `-- Services/
|   |       `-- ConsoleEmailSender.cs         <- implement this
|   |
|   `-- ShopFlow.Api/                     <- stub only for Stage 0
|       |-- ShopFlow.Api.csproj
|       |-- Program.cs
|       |-- Controllers/
|       |-- Middleware/
|       `-- Extensions/
|
`-- tests/
    |-- ShopFlow.Tests.Unit/
    |   |-- ShopFlow.Tests.Unit.csproj
    |   |-- Domain/
    |   |   |-- OrderTests.cs
    |   |   `-- ProductTests.cs
    |   |-- Application/
    |   |   |-- OrderServiceTests.cs      <- your main test file
    |   |   `-- PromoEngineTests.cs
    |   `-- Helpers/
    |       `-- DomainFactory.cs
    |
    `-- ShopFlow.Tests.Integration/       <- stub only for Stage 0
        |-- ShopFlow.Tests.Integration.csproj
        `-- Endpoints/
\end{codeblock}
\end{treebox}

\begin{dangerbox}
\textbf{Hard rules enforced at review:}\\[4pt]
\begin{itemize}[leftmargin=1.2em,itemsep=2pt]
  \item \texttt{ShopFlow.Domain.csproj} must have zero \texttt{<PackageReference>} entries
  \item \texttt{ShopFlow.Application} must not reference Infrastructure or Api
  \item \texttt{dotnet build} must pass with zero warnings
  \item \texttt{dotnet test} must pass with zero failures
\end{itemize}
\end{dangerbox}

% =============================================================================
\newpage
\section{The Broken Code}

A junior developer built the first working version of two service classes.
They compile. They run. They are wrong. Fix them.

% ─────────────────────────────────────────────────────────────────────────────
\subsection{File 1 --- src/ShopFlow.Domain/Entities/Order.cs}

\codelabel{Order.cs and OrderItem.cs as the junior left them}
\begin{codeframe}
\begin{codeblock}
public class Order
{
    public string          Id         { get; set; }
    public string          CustomerId { get; set; }
    public List<OrderItem> Items      { get; set; }
    public decimal         Total      { get; set; }
    public string          Status     { get; set; }  // "pending" "shipped" etc
    public DateTime        CreatedAt  { get; set; }
}

public class OrderItem
{
    public string  ProductId { get; set; }
    public string  Name      { get; set; }
    public decimal Price     { get; set; }
    public int     Quantity  { get; set; }
}
\end{codeblock}
\end{codeframe}

% ─────────────────────────────────────────────────────────────────────────────
\subsection{File 2 --- src/ShopFlow.Domain/Entities/Product.cs}

\codelabel{Product.cs as the junior left it}
\begin{codeframe}
\begin{codeblock}
public class Product
{
    public string  Id       { get; set; }
    public string  Name     { get; set; }
    public decimal Price    { get; set; }
    public int     Stock    { get; set; }
    public bool    IsActive { get; set; }

    public void Reserve(int quantity)
    {
        Stock -= quantity;   // no guard: stock can go negative
    }

    public void Restock(int quantity)
    {
        Stock += quantity;
    }
}
\end{codeblock}
\end{codeframe}

% ─────────────────────────────────────────────────────────────────────────────
\subsection{File 3 --- src/ShopFlow.Application/Services/OrderService.cs}

\codelabel{OrderService.cs --- every concern crammed into one class}
\begin{codeframe}
\begin{codeblock}
using System;
using System.Collections.Generic;
using System.Linq;
using System.Net.Mail;

public class OrderService
{
    private List<Order>   _orders   = new List<Order>();
    private List<Product> _products = new List<Product>();

    private string _smtpHost  = "smtp.shopflow.io";
    private int    _smtpPort  = 587;
    private string _emailFrom = "orders@shopflow.io";

    public string PlaceOrder(string customerId,
                             List<OrderItem> items,
                             string promoCode)
    {
        if (customerId == null || customerId == "")
            throw new Exception("Bad customer");
        if (items == null || items.Count == 0)
            throw new Exception("No items");

        // check stock inline
        foreach (var item in items)
        {
            var product = _products.FirstOrDefault(
                p => p.Id == item.ProductId);
            if (product == null)
                throw new Exception("Product not found: " + item.ProductId);
            if (product.Stock < item.Quantity)
                throw new Exception("Not enough stock for " + item.Name);
        }

        // calculate total
        decimal total = 0;
        foreach (var item in items)
            total += item.Price * item.Quantity;

        // apply promo
        if (promoCode == "SAVE10")
            total = total - (total * 0.10m);
        else if (promoCode == "SAVE20")
            total = total - (total * 0.20m);
        else if (promoCode == "HALFOFF")
            total = total / 2;
        // TODO: add more promo codes when marketing sends them

        // reserve stock inline
        foreach (var item in items)
        {
            var product = _products.First(p => p.Id == item.ProductId);
            product.Reserve(item.Quantity);
        }

        // create order
        var order        = new Order();
        order.Id         = Guid.NewGuid().ToString();
        order.CustomerId = customerId;
        order.Items      = items;
        order.Total      = total;
        order.Status     = "pending";
        order.CreatedAt  = DateTime.Now;
        _orders.Add(order);

        // send email directly
        try
        {
            var client  = new SmtpClient(_smtpHost, _smtpPort);
            var msg     = new MailMessage();
            msg.From    = new MailAddress(_emailFrom);
            msg.To.Add(customerId + "@customers.shopflow.io");
            msg.Subject = "Your order " + order.Id;
            msg.Body    = "Thanks! Total: " + total +
                          ". Items: " + items.Count;
            client.Send(msg);
        }
        catch (Exception e)
        {
            Console.WriteLine("Email failed: " + e.Message);
        }

        Console.WriteLine("[" + DateTime.Now + "] Order placed: "
            + order.Id + " customer=" + customerId + " total=" + total);

        return order.Id;
    }

    public Order GetOrder(string orderId)
    {
        foreach (var o in _orders)
            if (o.Id == orderId) return o;
        return null;
    }

    public bool CancelOrder(string orderId)
    {
        foreach (var o in _orders)
        {
            if (o.Id == orderId)
            {
                if (o.Status == "shipped")
                    return false;   // cannot cancel

                o.Status = "cancelled";

                // refund stock inline
                foreach (var item in o.Items)
                {
                    var product = _products.FirstOrDefault(
                        p => p.Id == item.ProductId);
                    if (product != null)
                        product.Restock(item.Quantity);
                }

                // copy-paste email block
                try
                {
                    var client  = new SmtpClient(_smtpHost, _smtpPort);
                    var msg     = new MailMessage();
                    msg.From    = new MailAddress(_emailFrom);
                    msg.To.Add(o.CustomerId + "@customers.shopflow.io");
                    msg.Subject = "Order " + orderId + " cancelled";
                    msg.Body    = "Your order has been cancelled.";
                    client.Send(msg);
                }
                catch (Exception e)
                {
                    Console.WriteLine("Email failed: " + e.Message);
                }
                return true;
            }
        }
        return false;   // not found -- same return value as "cannot cancel"
    }

    public List<Order> GetOrdersByCustomer(string customerId)
    {
        var result = new List<Order>();
        foreach (var o in _orders)
            if (o.CustomerId == customerId)
                result.Add(o);
        return result;
    }

    public decimal GetRevenueReport()
    {
        decimal revenue = 0;
        foreach (var o in _orders)
            if (o.Status != "cancelled")
                revenue += o.Total;
        return revenue;
    }
}
\end{codeblock}
\end{codeframe}

% ─────────────────────────────────────────────────────────────────────────────
\subsection{File 4 --- src/ShopFlow.Application/Services/InventoryService.cs}

\codelabel{InventoryService.cs --- also broken, separate state from OrderService}
\begin{codeframe}
\begin{codeblock}
using System;
using System.Collections.Generic;
using System.Linq;

public class InventoryService
{
    // completely separate list from the one in OrderService
    // stock changes here are invisible to OrderService and vice versa
    private List<Product> _products = new List<Product>();

    public void AddProduct(string id, string name,
                           decimal price, int initialStock)
    {
        if (string.IsNullOrEmpty(id))
            throw new Exception("Id required");
        if (price <= 0)
            throw new Exception("Bad price");

        _products.Add(new Product
        {
            Id       = id,
            Name     = name,
            Price    = price,
            Stock    = initialStock,
            IsActive = true
        });
    }

    public void AdjustStock(string productId, int delta)
    {
        var product = _products.FirstOrDefault(p => p.Id == productId);
        if (product == null)
            throw new Exception("Not found");

        product.Stock += delta;  // no guard: stock can go negative
    }

    public List<Product> GetLowStockProducts(int threshold)
    {
        // queries AND sends an email -- SRP violation
        var low = _products.Where(p => p.Stock < threshold).ToList();

        if (low.Count > 0)
        {
            try
            {
                var client = new System.Net.Mail.SmtpClient(
                    "smtp.shopflow.io", 587);
                var msg    = new System.Net.Mail.MailMessage();
                msg.From   = new System.Net.Mail.MailAddress(
                    "alerts@shopflow.io");
                msg.To.Add("warehouse@shopflow.io");
                msg.Subject = "Low stock alert";
                msg.Body    = low.Count + " products low.";
                client.Send(msg);
            }
            catch { }  // silently swallow all errors
        }

        return low;
    }

    public bool DeactivateProduct(string productId)
    {
        var product = _products.FirstOrDefault(p => p.Id == productId);
        if (product == null) return false;  // same false-for-two-reasons bug

        product.IsActive = false;
        return true;
    }
}
\end{codeblock}
\end{codeframe}

% =============================================================================
\newpage
\section{Problem Map --- Read After Forming Your Own List}

\begin{center}
\begin{tabularx}{\textwidth}{llX}
\toprule
\textbf{Principle} & \textbf{Broken?} & \textbf{Where} \\
\midrule
SRP & Yes & \texttt{OrderService} calculates totals, manages stock,
            sends emails, and logs --- four reasons to change \\
SRP & Yes & \texttt{InventoryService.GetLowStockProducts} queries
            and sends an alert email in the same method \\
OCP & Yes & Every new promo code requires editing \texttt{OrderService} \\
LSP & ---  & Not applicable yet \\
ISP & Yes & Callers needing only revenue reporting import
            the full \texttt{OrderService} with email and stock logic \\
DIP & Yes & Both services construct \texttt{SmtpClient} directly;
            nothing is injected \\
\midrule
State       & Critical & \texttt{OrderService} and \texttt{InventoryService}
                         each own a separate \texttt{List<Product>} so
                         stock changes in one are invisible to the other \\
Errors      & Weak     & \texttt{throw new Exception("Bad customer")} is
                         untyped and unactionable \\
Errors      & Weak     & \texttt{CancelOrder} returns \texttt{false} for
                         both not-found and already-shipped ---
                         caller cannot distinguish \\
Magic str.  & Yes      & \texttt{"pending"}, \texttt{"shipped"},
                         \texttt{"cancelled"} scattered across both files \\
Mutability  & Yes      & Every property on \texttt{Order}, \texttt{OrderItem},
                         and \texttt{Product} has a public setter \\
Guard       & Missing  & \texttt{Product.Reserve} allows stock to go negative \\
Dates       & Wrong    & \texttt{DateTime.Now} (local) instead of
                         \texttt{DateTimeOffset.UtcNow} \\
Duplication & Yes      & SMTP block copy-pasted verbatim in
                         \texttt{PlaceOrder}, \texttt{CancelOrder},
                         and \texttt{GetLowStockProducts} \\
\bottomrule
\end{tabularx}
\end{center}

% =============================================================================
\newpage
\section{Your Tasks}

\begin{taskbox}
Complete all five tasks and paste all code inline in the chat.
No GitHub repo needed for Stage 0.
\end{taskbox}

\subsection{Task 1 --- List Every Problem}

Before writing a single line of code, write your analysis as bullet points.
For each problem state: which principle is violated, which lines, and what the fix is.
This is what a senior developer does in a PR review.

\subsection{Task 2 --- Fix the Domain Entities}

Rewrite \texttt{Order}, \texttt{OrderItem}, and \texttt{Product} in the Domain project.

Requirements:
\begin{itemize}[leftmargin=1.5em,itemsep=3pt]
  \item No public setters --- use private setters and behaviour methods
  \item \texttt{OrderStatus} must be an enum, not a raw string
  \item \texttt{Order.Total} must be computed from items, not stored as a field
  \item \texttt{Product.Reserve(int)} must throw a typed \texttt{DomainException}
        if the quantity would make stock go negative
  \item All entities need an EF Core compatible private parameterless constructor
  \item Use \texttt{DateTimeOffset} not \texttt{DateTime}
  \item Create \texttt{DomainException} and \texttt{NotFoundException}
        under \texttt{ShopFlow.Domain/Exceptions/}
\end{itemize}

\subsection{Task 3 --- Refactor the Application Layer}

Rewrite both services so they satisfy SOLID. You must:

\begin{itemize}[leftmargin=1.5em,itemsep=3pt]
  \item Create interfaces \texttt{IOrderRepository}, \texttt{IProductRepository},
        and \texttt{IEmailSender} in \texttt{ShopFlow.Application/Interfaces/}
  \item Create a \texttt{PromoEngine} class that is open for extension ---
        adding a new promo code must not require editing \texttt{OrderService}
  \item \texttt{OrderService} must depend only on those interfaces via
        constructor injection, never on concretions
  \item \texttt{InventoryService} must use \texttt{IProductRepository}
        and \texttt{IEmailSender}, never \texttt{SmtpClient} directly
  \item \texttt{CancelOrder} must not return \texttt{false} for two different
        reasons --- fix the ambiguity however you see fit
  \item Implement \texttt{InMemoryOrderRepository} and
        \texttt{InMemoryProductRepository} in Infrastructure
  \item Implement \texttt{ConsoleEmailSender} in Infrastructure
        (writes to console --- no real SMTP needed)
\end{itemize}

\subsection{Task 4 --- Write Unit Tests}

Write at least \textbf{10 unit tests} in \texttt{OrderServiceTests.cs}
using xUnit and Moq. Cover:

\begin{itemize}[leftmargin=1.5em,itemsep=3pt]
  \item \texttt{PlaceOrder} happy path --- order saved, id returned
  \item \texttt{PlaceOrder} empty \texttt{customerId} --- correct exception
  \item \texttt{PlaceOrder} empty items list --- correct exception
  \item \texttt{PlaceOrder} product not found --- correct exception
  \item \texttt{PlaceOrder} insufficient stock --- correct exception
  \item \texttt{PlaceOrder} with \texttt{SAVE10} promo --- total discounted correctly
  \item \texttt{PlaceOrder} with unknown promo --- no discount applied
  \item \texttt{PlaceOrder} email failure --- order is still saved
  \item \texttt{CancelOrder} on a pending order --- succeeds
  \item \texttt{CancelOrder} on a shipped order --- fails with the correct signal
\end{itemize}

Write at least \textbf{4 domain tests} in \texttt{OrderTests.cs}
and \texttt{ProductTests.cs}.

\subsection{Task 5 --- Answer in Prose}

\begin{notebox}
\textbf{Answer both questions before you submit.}\\[6pt]
\textbf{Q1.} The original code has two separate \texttt{List<Product>} ---
one in \texttt{OrderService} and one in \texttt{InventoryService}.
Stock reserved during \texttt{PlaceOrder} is invisible to
\texttt{InventoryService}. How does your refactoring fix this,
and which design principle does it invoke?\\[8pt]
\textbf{Q2.} You chose a specific strategy for the promo code extensibility problem.
What pattern did you use, why is it open for extension,
and what is its trade-off?
\end{notebox}

% =============================================================================
\section{Acceptance Checklist}

Run these yourself before pasting your solution:

\begin{greenbox}
\begin{itemize}[leftmargin=1.2em,itemsep=4pt]
  \item[$\square$] \texttt{dotnet build} --- zero errors, zero warnings
  \item[$\square$] \texttt{dotnet test} --- all tests green
  \item[$\square$] \texttt{ShopFlow.Domain.csproj} has zero \texttt{<PackageReference>} entries
  \item[$\square$] No \texttt{SmtpClient} anywhere outside the Infrastructure project
  \item[$\square$] No magic strings \texttt{"pending"} / \texttt{"shipped"} / \texttt{"cancelled"}
  \item[$\square$] No public setters on any domain entity
  \item[$\square$] \texttt{Product.Reserve} throws on negative stock
  \item[$\square$] \texttt{CancelOrder} distinguishes not-found from cannot-cancel
  \item[$\square$] All service methods accept a \texttt{CancellationToken} parameter
  \item[$\square$] At least 14 passing unit tests
\end{itemize}
\end{greenbox}

\end{document}
